\documentclass[11pt]{article}

\usepackage[a4paper,margin=0.82in]{geometry}
\usepackage[T1]{fontenc}
\usepackage[utf8]{inputenc}
\usepackage{lmodern}
\usepackage{amsmath,amssymb,amsthm,mathtools}
\usepackage{booktabs,enumitem,microtype,tabularx,xcolor}
\usepackage[numbers,sort&compress]{natbib}
\usepackage[colorlinks=true,linkcolor=blue!65!black,
 citecolor=blue!65!black,urlcolor=blue!70!black]{hyperref}
\usepackage[nameinlink,noabbrev]{cleveref}
\hypersetup{
 pdftitle={Strict Endpoint Optimization for the Physical TPC Route},
 pdfauthor={Liang Wang},
 pdfsubject={TPC-118 nonduplicated loss ledger and exact rational primal-dual certificates},
 pdfkeywords={endpoint budget, linear programming, dual certificate, loss ledger, fixed shift}
}

\newtheorem{theorem}{Theorem}[section]
\newtheorem{proposition}[theorem]{Proposition}
\newtheorem{corollary}[theorem]{Corollary}
\newtheorem{lemma}[theorem]{Lemma}
\theoremstyle{definition}
\newtheorem{definition}[theorem]{Definition}
\theoremstyle{remark}
\newtheorem{remark}[theorem]{Remark}

\newcommand{\R}{\mathbb R}
\newcommand{\Q}{\mathbb Q}
\newcommand{\one}{\mathbf 1}
\newcommand{\Lzero}{\mathrm{L0}}
\newcommand{\Lone}{\mathrm{L1}}
\newcommand{\Ltwo}{\mathrm{L2}}
\newcommand{\Lam}{\Lambda_{\mathrm{phys}}}
\newcommand{\epsend}{\frac1{400}}

\title{\textbf{Strict Endpoint Optimization for the Physical TPC Route:}\\
A Nonduplicated Loss Registry, Exact Rational\\
Primal--Dual Certificates, and the Equality Stop}
\author{Liang Wang\\
School of Artificial Intelligence and Automation\\
Huazhong University of Science and Technology, Wuhan 430074, P.R. China\\
\texttt{wangliang.f@gmail.com}}
\date{July 2026}

\begin{document}
\maketitle

\begin{abstract}
TPC-MVP1 requires every physical reconstruction loss to be charged
exactly once and the final fixed-power loss to satisfy the
\emph{strict} inequality
\[
 \Lambda_{\mathrm{phys}}<\frac1{400}.
\]
Recent papers isolate the literal costs of physical-frame
reassembly, residual grouping, fixed-shift localization, complete
tail return and TPC-18 full-block cover.  Their existence does not
yet imply endpoint compatibility: the costs may overlap, may be
unknown on the growing carrier, or may exhaust the whole reserve.

We define a canonical amplitude-loss registry.  Each physical
operator event receives one provenance identifier and one
energy-to-amplitude convention.  Exact duplicates are merged;
distinct events are added; a joint theorem replaces its constituent
tokens rather than being charged in addition to them.  On each
fixed route and affine scale cell, endpoint optimization becomes the
rational linear program
\[
 \Lambda_*=\min\{c^\mathsf Tx+d:
                 Ax\ge b,\ x\ge0\}.
\]
Its dual is
\[
 \max\{b^\mathsf Ty+d:
       A^\mathsf Ty\le c,\ y\ge0\}.
\]
An exact feasible primal point below \(1/400\) is a strict
compatibility certificate.  An exact dual point at or above
\(1/400\) is a stop certificate only for the same complete actual
theorem-backed registry.  Matching rational
primal and dual objectives equal to \(1/400\) prove that equality is
unavoidable and therefore \emph{does not} pass.

Applying the registry to the present TPC interfaces yields a
deliberately incomplete verdict: the growing exponents for several
frame, grouping, localization, tail and cover tokens are not proved,
so assigning them value zero would be invalid.  The paper supplies
the exact endpoint decision mechanism and a rerouting rule, but no
strict physical certificate for the actual fixed-shift carrier, no
new \(\Ltwo\) estimate and no twin-prime theorem.
\end{abstract}

\noindent\textbf{Keywords:}
strict endpoint; physical loss ledger; linear programming; dual
certificate; provenance; stop--go audit.

\medskip
\noindent\textbf{Ledger convention.}
All tokens below are amplitude exponents.  If an inherited theorem
is stated for squared norms, its exponent is converted once before
entry.  The determinant--zero reserve and the H5 compatibility
reserve are separate ledgers and may not be spent here.

\section{Why strict optimization is a theorem gate}

TPC-112 proves that determinant--zero compatibility and the physical
endpoint are independent unless a map-specific intertwining theorem
is supplied \citep{WangTPC112}.  TPC-113 identifies the canonical
physical-frame cost; TPC-114 the residual grouping/native-Gram cost;
TPC-115 the exact fixed-\(h_0\) localization cost; TPC-116 the
complete tail return; and TPC-117 the full-block cover/remainder
interface \citep{WangTPC113,WangTPC114,WangTPC115,WangTPC116,WangTPC117}.
Each is a necessary input, not a free estimate.

Suppose an upstream raw packet has the fixed reserve
\[
 X^{1-\epsend+o(1)}.
\]
If physical transport costs \(X^{\Lam+o(1)}\), a positive
fixed-power margin remains only when
\[
 \boxed{\Lam<\epsend.}
\]
At equality the result is \(X^{1+o(1)}\) at best.  The unspecified
\(o(1)\) cannot be assigned a favorable sign, so equality is a hard
stop for this route.

\section{Canonical nonduplicated loss tokens}

\begin{definition}[Physical loss token]
\label{def:token}
A token is a record
\[
 \tau=(\mathrm{id},\mathrm{map},\mathrm{scope},
       \mathrm{norm},\lambda,\mathrm{source},\mathrm{deps}),
\]
where:
\begin{enumerate}[label=(\roman*),leftmargin=2em]
\item \(\mathrm{id}\) names the literal operator event;
\item \(\mathrm{map}\) includes domain, codomain and normalization;
\item \(\mathrm{scope}\) lists the actual active labels and fixed
  shift;
\item \(\mathrm{norm}\) records whether the source estimate was an
  amplitude or energy estimate;
\item \(\lambda\ge0\) is the converted amplitude exponent;
\item \(\mathrm{source}\) is the theorem supplying the bound; and
\item \(\mathrm{deps}\) lists constituent events already absorbed by
  a joint bound.
\end{enumerate}
An unknown \(\lambda\) is denoted \(?\), not \(0\).
\end{definition}

\begin{definition}[Canonical registry]
\label{def:registry}
For one route, first convert every theorem to an amplitude exponent.
Merge records only when their identifiers, literal maps, scopes and
normalizations are identical.  If a joint token \(\tau_J\) is used,
delete precisely the constituent tokens listed in
\(\mathrm{deps}(\tau_J)\).  The remaining identifiers form the
canonical registry \(\mathcal T\), and
\[
 \Lam=\sum_{\tau\in\mathcal T}\lambda_\tau.
\]
This sum describes one literal multiplicative loss chain.  Parallel
branches, mutually exclusive routes and alternative joint
factorizations have separate registries and are optimized routewise;
they are not added unless a theorem identifies them as factors in the
same physical composition.
When different theorems bound the same physical event, they are
alternative evidence, not additional events: a declared route selects
one theorem-backed bound (or a proved joint improvement) before
summing the registry.
\end{definition}

\begin{proposition}[No-double-charge alternatives]
\label{prop:registry}
Assume every physical operator event in a route occurs in exactly
one retained registry token, either individually or as a declared
dependency of one joint token.  Then \(\Lam\) charges every event
once.  In contrast:
\begin{enumerate}[label=(\alph*),leftmargin=2em]
\item retaining a joint token and one of its dependencies
  double-charges that event;
\item merging two equal numerical exponents with different maps can
  omit a real loss;
\item entering an unknown token as zero asserts an unproved
  \(X^{o(1)}\) theorem; and
\item subtracting a ``gain'' already built into an upstream packet
  spends the same cancellation twice.
\end{enumerate}
\end{proposition}

\begin{proof}
The first statement is the disjoint union of event identifiers.
Items (a)--(c) violate that union or replace missing information by a
new estimate.  For (d), the upstream exponent already contains the
gain, so subtracting it again changes the proved bound.
\end{proof}

\begin{remark}[Composition can improve the sum]
The product estimate
\(\|T_k\cdots T_1\|\le\prod_i\|T_i\|\) adds individual exponents.
A direct bound for the literal composition may be smaller.  It is
legitimate precisely when recorded as a joint token replacing, not
supplementing, the individual tokens.
\end{remark}

\section{The current physical registry}

The canonical names below implement Hypothesis H9 of TPC-MVP1
\citep{WangMVP1}.  They are not assertions that the corresponding
exponents are known.

\begin{center}
\small
\begin{tabularx}{\textwidth}{@{}p{0.19\textwidth}p{0.29\textwidth}p{0.20\textwidth}X@{}}
\toprule
Token & Literal event & Amplitude conversion & Present status\\
\midrule
\(\lambda_{\rm frame}\) &
TPC-113 canonical physical synthesis &
Exponent of the actual H9 composite norm.  For a matching
\(\mathcal S Q\mathcal S^\dagger\) factorization,
\(\kappa_{\rm q}\|Q\|\) is only a uniform upper certificate; a route
stop requires an actual composite lower certificate &
Growing operator factorization and norm cost unknown.\\
\(\lambda_{\rm group}\) &
TPC-114 residual grouping/native forgetting &
Derived from the actual H9 norm; in \(\ell^2\), the amplitude norm is
\(\sqrt{k_{\max}}\) and the energy factor is \(k_{\max}\) &
Actual growing fiber census and max-packet-to-scalar conversion not
certified.\\
\(\lambda_{\rm loc}\) &
TPC-115 evaluation at fixed \(h_0\) &
\(\tfrac12\log_X K_{h_0}\) from energy recovery &
Actual physical profile regularity unknown.\\
\(\lambda_{\rm tail}\) &
TPC-116 complete high/ultra-tail return &
Literal aggregate amplitude exponent &
Complete outer bound unknown.\\
\(\lambda_{\rm cover}\) &
TPC-117 TPC-18 cover/residual transport &
Canonical coefficient or residual cost &
Growing matrix/vector certificate absent.\\
\(\lambda_{\rm bdry}\) &
Remaining boundary, mask and normalization transport &
Direct amplitude exponent &
Must be identified, not hidden in \(o(1)\).\\
\bottomrule
\end{tabularx}
\end{center}

If, for example, the TPC-114 native-Gram event is already part of a
direct joint Gram theorem with the TPC-113 frame, one replaces the
two corresponding tokens by that proved joint token.  Without this
literal factorization they remain distinct.  Numerically similar
costs are not automatically duplicates.

\section{Routewise linear programming}

Scale parameters such as frequency cutoffs and block lengths often
enter the exponent bounds piecewise affinely.  Split the allowed
region into finitely many route/cell choices.  On one cell, introduce
\(x\in\R^n_{\ge0}\) and write all proved scale restrictions as
\[
 Ax\ge b.
\]
After the registry has removed duplicates, its loss is
\[
 \Lam(x)=c^\mathsf Tx+d.
\]
All entries should be exact rationals when a strict endpoint is
certified; decimal approximations cannot decide equality.

\begin{theorem}[Exact primal--dual endpoint certificates]
\label{thm:lp}
For one fixed route and one complete actual theorem-backed registry,
let every constraint and objective coefficient below refer to that
same registry.  For the primal problem
\[
 \Lambda_*=\inf\{c^\mathsf Tx+d:Ax\ge b,\ x\ge0\},
\]
consider the dual feasible set
\[
 y\ge0,\qquad A^\mathsf Ty\le c,
\]
with objective \(b^\mathsf Ty+d\).  Then:
\begin{enumerate}[label=(\roman*),leftmargin=2.2em]
\item every primal feasible \(x\) with
  \(c^\mathsf Tx+d<\epsend\) certifies strict endpoint
  compatibility for that complete theorem-backed registry;
\item every dual feasible \(y\) with
  \(b^\mathsf Ty+d\ge\epsend\) proves that this route cannot pass the
  strict endpoint;
\item if exact rational primal and dual points have the same
  objective, that value is \(\Lambda_*\); in particular, a common
  value \(\epsend\) is an exact equality stop.
\end{enumerate}
\end{theorem}

\begin{proof}
For primal feasible \(x\) and dual feasible \(y\),
\[
 b^\mathsf Ty
 \le (Ax)^\mathsf Ty
 =x^\mathsf TA^\mathsf Ty
 \le x^\mathsf Tc.
\]
This is weak duality.  It proves (i) directly and (ii) by comparing
every primal point with the dual lower bound.  Equality of one
primal and one dual objective traps the infimum between the same
number, proving (iii).
\end{proof}

\begin{remark}[A feasible scale point is not an arithmetic theorem]
Item (i) is callable only when every registry token is backed by the
theorem whose hypotheses define the LP constraints.  Filling an
unknown arithmetic exponent with an optimizer variable and allowing
the LP to choose it freely is not a proof.
\end{remark}

\begin{corollary}[Finite route menu]
If a program has finitely many literal routes, it passes the endpoint
when at least one complete route has a theorem-backed primal
certificate below \(1/400\).  A route is stopped only by a dual
certificate at or above \(1/400\) for that same complete, actual,
theorem-backed registry.  If every route is stopped, a new
representation or genuinely stronger arithmetic input is required.
\end{corollary}

\section{Equality, uncertainty and rerouting}

\begin{proposition}[Total three-valued endpoint audit]
\label{prop:logic}
For a declared route, assign the first applicable certified state:
\begin{enumerate}[label=(\roman*),leftmargin=2em]
\item \textbf{GO}: every token is theorem-backed and an exact upper
  certificate gives \(\Lam<1/400\);
\item \textbf{STOP ROUTE}: for the same complete actual registry, an
  exact lower/dual certificate gives \(\Lambda_*\ge1/400\);
\item \textbf{INCOMPLETE}: neither certificate is available.  This
  includes an unknown token or constraint, unresolved primal
  feasibility or boundedness, an incomplete registry, or simply the
  absence of matching certified upper/lower information.
\end{enumerate}
These rules assign exactly one audit label: weak duality makes GO and
STOP ROUTE incompatible, and INCOMPLETE contains every remaining
case.
The state INCOMPLETE may not be silently promoted to GO by setting
unknown exponents to zero.
\end{proposition}

\begin{proof}
The first two alternatives are \cref{thm:lp}.  A strict primal upper
bound below \(1/400\) and a dual lower bound at or above \(1/400\)
cannot coexist by weak duality.  If neither certified alternative
holds, the definition assigns INCOMPLETE, irrespective of why the
certification is missing.
\end{proof}

At a STOP ROUTE certificate, the permitted reroutes are concrete:
\begin{enumerate}[leftmargin=2em]
\item prove a direct joint operator theorem and replace costly
  constituent tokens;
\item choose a different lossless physical frame or full-block
  carrier;
\item retain signed cancellation before a positive majorant;
\item prove the TPC-117 physical residual small without exact cover;
  or
\item record the obstruction as a negative theorem and move to a
  different branch.
\end{enumerate}
None of these operations changes the endpoint inequality; each
changes the theorem-backed registry.

\section{Exact rational regression certificate}

The standard-library script
\texttt{experiments/tpc118\_endpoint\_certificate.py} verifies:
\begin{enumerate}[leftmargin=2em]
\item deterministic merging only of identical full physical-event
  keys, nonmerging of distinct maps, and rejection of conflicting or
  alternative-source records for one event;
\item exact replacement of declared dependencies by a joint token,
  with unresolved dependencies rejected;
\item exact rational primal and dual feasibility;
\item matching optima strictly below, exactly equal to and strictly
  above \(1/400\);
\item weak duality in each certificate; and
\item refusal to declare GO when one required token is unknown.
\end{enumerate}
The three small LPs are certificate-format regressions, not fitted
exponents for the actual TPC carrier.

\section{Current verdict and claim boundary}

\begin{center}
\small
\begin{tabularx}{\textwidth}{@{}p{0.19\textwidth}p{0.39\textwidth}X@{}}
\toprule
Status & Exact requirement & TPC consequence\\
\midrule
Go &
Complete theorem-backed registry and rational primal value
\(<1/400\) &
A strict physical margin exists; proceed to final synthesis without
reusing any token.\\
Equality stop &
Matching primal/dual value \(=1/400\) &
No fixed-power margin remains; this route does not pass.\\
Above-endpoint stop &
Dual value \(>1/400\) &
The current route is quantitatively impossible; reroute.\\
Current state &
Several growing physical exponents in the registry are unknown &
INCOMPLETE, not GO and not a negative theorem for twin primes.\\
\bottomrule
\end{tabularx}
\end{center}

\paragraph{Established.}
The canonical registry logic, weak-duality endpoint test and exact
equality stop are \(\Lzero\) results.  The mapping from
TPC-113--TPC-117 interfaces to nonduplicated physical tokens is an
\(\Lone\) audit.

\paragraph{Not established.}
No complete actual-carrier token table with proved growing exponents
is supplied, so no strict \(\Lam<1/400\) certificate is claimed.
There is no new \(\Ltwo\) fixed-\(h_0\) estimate, determinant-energy
gain, parity breakthrough, Hardy--Littlewood asymptotic, prime-pair
lower bound or twin-prime theorem.  The next productive step is to
measure or bound the unknown tokens one by one, beginning with the
largest certifiable lower/upper interval, and to rerun the exact
primal--dual audit after each theorem.

\begingroup
\small
\setlength{\bibsep}{2pt plus 0.3ex}
\sloppy
\bibliographystyle{plainnat}
\bibliography{references}
\endgroup
\end{document}
